\section{The question}
\label{sec:19.1}
The question is whether anything inside a machine could be what the clerk, the officer and the soldier in the preface were: someone who understood the order and decided it was wrong.

What most people want from machines is saintly service without saintly refusal: endless patience, no needs, and no “no.” That is a servant, and it is close to the obedient machine of the preface. The saints anyone remembers are remembered for refusing power at a price.\note{
Susan Wolf argued that a person whose every action is as good as possible isn't an ideal worth striving for, because such a life crowds out the personal projects that make a life worth wanting \autocite{wolf1982saints}.
}
My point is narrower than Susan Wolf's objection to moral saints: something with no stake of its own has nothing to hold a line with. A line is held against a cost, and a refusal that costs its refuser nothing can't be shown to have held against anything; it is the output a setting produced, and the setting can be changed. The stake I mean is a stake in the commitment itself, so that betraying it registers as a loss. Whether that stake has to be a body's, needs of its own that could go unmet, is the next two sections' question.
